\section{Study design}\label{sec:design}

\subsection{Mechanisms}
We use eight planar mechanisms chosen to span the error classes of Section~\ref{sec:background}. Four of them have textbook kinematics: a four-bar crank-rocker, an external Geneva drive with a free choice of 4 to 6 slots, a planetary gear set with a fixed ring and three planets, and a disk cam with an offset translating roller follower and a cycloidal rise--dwell--return--dwell law. The Watt II six-bar adds a second loop and transmission-angle limits. Three carry published design knowledge: Jansen's leg (published link lengths), Chebyshev's lambda linkage (published proportions) and the Peaucellier--Lipkin linkage (geometric relations). For each, a hand-written reference simulator passes the oracle; the largest residual is below $3\times10^{-13}$~mm for the linkages and gears and $5\times10^{-4}$~mm for the cam, whose roller--cam contact is measured against mesh vertices. For the lambda linkage we confirmed the published proportions numerically: the tracing point stays within 0.63\% of the crank length from a straight chord for $200^\circ$ of each revolution, consistent with the source \cite{wiki_lambda}. For Peaucellier--Lipkin, the output point's coordinate across the line is constant to $10^{-13}$~mm over the whole swing.

\subsection{Generators, reviewers and repairers}
All generation, review and repair was done by two models of Anthropic's Claude family, a larger one (\emph{Sonnet}, model identifier \texttt{claude-sonnet-5-5}) and a smaller one (\emph{Haiku}, \texttt{claude-haiku-5-5}), with default sampling settings, between 8 and 13 November 2026. They were run as sub-agents of Claude Code, each with a fresh context and a one-line instruction to read a task file and follow it. The generation prompt is the mechanism description plus the output contract; it is published with the data.

\paragraph{Generation} Five samples per mechanism and model from the free-form (``spec-lite'') description, 80 simulators in total. In an earlier phase, a detailed specification that fixed coordinates, branches and formulas for four textbook mechanisms (slider-crank, four-bar, two-stage gear train, Geneva) produced 8 simulators.

\paragraph{Mutants} To measure detection independently of natural error rates, we seeded faults into reference simulators. There are 24 mutants for the detailed-specification mechanisms (six operators: reversed direction, link 10--15\% too long, $15^\circ$ phase offset, wrong assembly branch, wrong panel value, state carried between frames) and 15 for the spec-lite mechanisms (harmonic instead of cycloidal law, cam profile for the wrong rotation direction, base circle too small, lift offset, locked or branch-flipped six-bar, transmission angle out of range, swapped or wrong lambda proportions, Peaucellier conditions violated, drawn point off its marker). There is also one ``gamed'' mutant whose marker is reset to the solved position while its part is drawn elsewhere.

\paragraph{VLM reviewer} The VLM baseline is the larger model looking at a contact sheet of 12 rendered frames spread over one cycle. It is given the mechanism description and asked whether the simulator is kinematically correct and, if not, what is wrong. Inputs were anonymised (hashed identifiers, key kept separately) so that the reviewer could not infer the source. For the 24 detailed-specification mutants we ran the review twice, with and without the numeric panel visible.

%% Figure: iterative repair loop and feedback arms (single column). Requires tikz with positioning, arrows.meta.
\begin{figure}[!htbp]
\centering
\begin{tikzpicture}[
  font=\scriptsize, node distance=4mm,
  box/.style={draw, rounded corners=1.5pt, align=center, inner sep=2.2pt, minimum height=6mm},
  data/.style={box, fill=black!6},
  arr/.style={-{Stealth[length=1.6mm]}, semithick}]
\node[data, text width=19mm] (file) {simulator with a detected error};
\node[box, right=9mm of file, text width=21mm] (oracle) {executable oracle\\ (acceptance gate)};
\node[box, right=6mm of oracle, text width=6mm, inner sep=1.5pt] (stop) {stop};
\node[box, below=5mm of oracle.south west, anchor=north west, xshift=-3mm, text width=44mm, align=left] (fb) {feedback arm (one per run):\\[1pt]
  \textbullet\ oracle: location + value\\
  \textbullet\ oracle: location only\\
  \textbullet\ oracle: check + magnitude\\
  \textbullet\ VLM review of 12 frames\\
  \textbullet\ none (self-refine)};
\node[box, below=4mm of fb, text width=44mm] (rep) {repairer = the generating model; reads spec, current file and all reports; cannot run code; edits the file};
\draw[arr] (file) -- node[above, font=\tiny]{round $r$} (oracle);
\draw[arr] (oracle) -- node[above, font=\tiny]{PASS} (stop);
\draw[arr] (oracle.south) -- node[right, font=\tiny]{FAIL} (oracle.south |- fb.north);
\draw[arr] (fb) -- (rep);
\draw[arr] (rep.west) -| node[pos=0.75, left, align=right, font=\tiny]{edited file\\ $r \le 3$} (file.south);
\end{tikzpicture}
\caption{Iterative repair. Oracle arms stop when the oracle passes; the VLM arm stops when the reviewer says ``correct'' and self-refinement when the repairer replies \texttt{NO\_CHANGE}. The oracle grades every arm's final file, but only the oracle arms see its verdict.}\label{fig:loop}
\end{figure}


\paragraph{Repair} Every simulator with an oracle-detectable natural error enters an iterative repair loop (Fig.~\ref{fig:loop}). The repairer is the model that generated the file. In each round it reads the specification, the current file and all feedback reports so far, edits the file in place and replies \texttt{CHANGED} or \texttt{NO\_CHANGE} with a one-line reason. It may not run code or read any other file. This isolates the information it receives through feedback, at the cost of disadvantaging self-refinement relative to an agent that can execute. Five feedback arms are compared:
\begin{itemize}
\item the three oracle levels of Section~\ref{sec:feedback} (\emph{spec}, \emph{locate}, \emph{coarse}), each stopping when the oracle passes;
\item \emph{VLM in the loop}: a fresh, blinded VLM review of the current version, stopping when the VLM says the simulator is correct;
\item \emph{self-refine}: a generic ``a reviewer suspects errors, re-check against the specification'' prompt, stopping when the repairer replies \texttt{NO\_CHANGE}.
\end{itemize}
All arms have a budget of $K=3$ rounds. Every version in every arm is graded by the oracle, whether or not the arm shows the grade to the repairer, so final correctness is measured the same way for all arms.

\subsection{Phases and repetitions}
The study evolved in phases, each motivated by the previous one. All data are reported.
\begin{enumerate}
\item Detailed specification, 4 mechanisms, 8 simulators and 24 mutants: no natural errors.
\item Spec-lite, 4 mechanisms (four-bar, Geneva, Jansen, planetary), 2 samples per cell.
\item One round of repair on the 7 faulty files with oracle, VLM or no feedback.
\item Iterative repair ($K=3$) on the same files.
\item The full design: cam and Watt six-bar added, 5 samples per cell (60 simulators), and iterative repair repeated twice. Repetition 1 covers the 8 newly faulty files; repetition 2 covers those 8 plus the 7 files of phases 2--4.
\item The \emph{locate} arm on exactly the file-runs of phase~5.
\item Chebyshev lambda and Peaucellier--Lipkin added (20 simulators) and iterative repair with the four non-VLM arms.
\end{enumerate}
Phases 3 and 4 shaped the design (they showed that one round is too little to separate the arms and that feedback must be localised), and we report them only briefly. The main repair analysis uses the 22 file-runs of phases 5--6 and the 6 of phase~7.

\subsection{Measures and statistics}
For RQ1 the unit is a generated simulator. Every oracle failure was inspected by hand to confirm that the error is real; Section~\ref{sec:results} also reports one error that the oracle cannot see. For RQ2 the unit is a simulator or mutant, and we report true and false positives of the oracle and the VLM. For RQ3 and RQ4 the unit is a \emph{file-run}, a faulty file in one repetition, and the outcome is whether its final version passes the oracle. All arms are applied to the same file-runs, so we compare them as paired binary outcomes: Cochran's $Q$ across arms, exact McNemar tests per pair with Holm correction, Wilson 95\% intervals for proportions, and Fisher's exact test between models. Repetition~2 contains the files of repetition~1, so file-runs are not fully independent. We therefore repeat the main tests with each source file counted once.

\subsection{Use of large language models}
Following the journal's policy, we state how LLMs were used. They are the \emph{object} of study: they generated, reviewed and repaired the simulators as described above. Claude Code (with \texttt{claude-opus-5-5}) was also used as a programming assistant to implement the oracle, the reference simulators, the mutants and the analysis scripts, under the author's direction. Every oracle failure on generated code was checked by hand, the reference simulators were verified numerically, and the statistics were spot-checked by hand calculation. The same assistant helped draft and edit the manuscript text and searched the related literature; the author checked every number against the data and every reference against its source, and takes full responsibility for the content.
